%------------------------------------------------------------------------------%
\documentclass[fleqn,utf8,aspectratio=169,12pt]{beamer}

\usepackage{gaal}

\title{Classificação de Sistemas Lineares}

%------------------------------------------------------------------------------%
\begin{document}

\addtitle

%------------------------------------------------------------------------------%
\begin{frame}{Tipos de Sistemas Lineares}
  \onslide<+->{Um sistema pode ser classificado com relação ao número de soluções}
  \begin{itemize}
    \setlength{\itemsep}{1em}
    \item \onslide<+->{Sistema possível e determinado} \\
          \onslide<+->{possui \emph{uma única} solução}

    \item \onslide<+->{Sistema possível e indeterminado} \\
          \onslide<+->{possui \emph{infinitas} soluções}

    \item \onslide<+->{Sistema impossível} \\
          \onslide<+->{\emph{não possui} soluções}
  \end{itemize}
\end{frame}

%------------------------------------------------------------------------------%
\begin{frame}{Identificando o Tipo de Sistema Linear}
  Podemos identificar o tipo de sistema diretamente da sua forma escalonada
\end{frame}

%------------------------------------------------------------------------------%
%------------------------------------------------------------------------------%
\begin{question}
  Sistema possível e determinado
  \[
    \begin{amatrix}{3}
      \emph{1} & a        & b        & c \\
      0        & \emph{1} & d        & e \\
      0        & 0        & \emph{1} & f
    \end{amatrix}
  \]

  \pause
  Possui uma única solução
\end{question}
%------------------------------------------------------------------------------%

%------------------------------------------------------------------------------%
\begin{question}
  Sistema possível e indeterminado
  \[
    \begin{amatrix}{3}
      \emph{1} & a        & b & c \\
      0        & \emph{1} & d & e \\
      0        & 0        & 0 & 0
    \end{amatrix}
  \]

  \pause
  Possui infinitas soluções

  \pause
  $x$ e $y$ ficam em função de $z$ que pode assumir qualquer valor
\end{question}

%------------------------------------------------------------------------------%

%------------------------------------------------------------------------------%
\begin{question}
  Sistema impossível

  Se uma linha tiver a forma
  \[
    \begin{amatrix}{3}
      \emph{1} & a        & b & c \\
      0        & \emph{1} & d & e \\
      0        & 0        & 0 & f
    \end{amatrix}
    \qquad f \neq 0
  \]

  \pause
  Não possui solução

  \pause
  Não importa os valores de $x$, $y$ ou $z$
  \[
    0x + 0y + 0y = 0 \neq f
  \]
\end{question}

%------------------------------------------------------------------------------%

%------------------------------------------------------------------------------%
\begin{question}
Use escalonamento para analisar o sistema
\[
  \systeme{
     x +  y = 3,
    2x + 2y = 8
  }
\]
\end{question}

%------------------------------------------------------------------------------%
\begin{resolution}{Escalonamento}
\begin{columns}
\column{0.4\linewidth}
  \[\;\;\,
    \begin{amatrix}{2}
      1 & 1 & 3 \\
      2 & 2 & 8
    \end{amatrix}
  \]
  \pause
\column{0.6\linewidth}
  Eliminar os elementos abaixo do primeiro pivô
  \\ \pause
  \(
    L_2 \leftarrow L_2 - 2 L_1
  \)
\end{columns}
\vspace{-\baselineskip}
\[
  \pause
  L'_2
  \;=\;
  L_2 - 2 L_1
  \pause
  \;=\;
  \begin{amatrix}{2}
    2 & 2 & 8
  \end{amatrix}
  - 2
  \begin{amatrix}{2}
    1 & 1 & 3
  \end{amatrix}
  \pause
  \;=\;
  \begin{amatrix}{2}
    0 & 0 & 2
  \end{amatrix}
\]

\pause
\[
\begin{amatrix}{2}
  1 & 1 & 3 \\
  0 & 0 & 2
\end{amatrix}
\]
\end{resolution}

%------------------------------------------------------------------------------%
\begin{resolution}{Solução}
A segunda linha representa
\[
  0 = 2
\]
\pause
o que é uma contradição

\pause
\bigskip
Portanto, o sistema é \emph{incompatível} e não possui solução
\end{resolution}

%------------------------------------------------------------------------------%
\endinput

%------------------------------------------------------------------------------%
\begin{question}
Use escalonamento para analisar o sistema
\[
  \systeme{
   x +  y = 4,
  2x + 2y = 8
  }
\]
\end{question}

%------------------------------------------------------------------------------%
\begin{resolution}{Escalonamento}
\begin{columns}
\column{0.4\linewidth}
  \[\;\;\,
    \left[
    \begin{array}{cc|c}
      1 & 1 & 4 \\
      2 & 2 & 8
    \end{array}
    \right]
  \]
  \pause
\column{0.6\linewidth}
  Aplicamos
  \\ \pause
  \(
    L_2 \leftarrow L_2 - 2 L_1
  \)
\end{columns}
\vspace{-\baselineskip}
\[
  \pause
  L'_2
  \;=\;
  L_2 - 2 L_1
  \pause
  \;=\;
  \left[
  \begin{array}{cc|c}
    2 & 2 & 8
  \end{array}
  \right]
  - 2
  \left[
  \begin{array}{cc|c}
    1 & 1 & 4
  \end{array}
  \right]
  \pause
  \;=\;
  \left[
  \begin{array}{cc|c}
    0 & 0 & 0
  \end{array}
  \right]
\]

\pause
\[
\left[
\begin{array}{cc|c}
  1 & 1 & 4 \\
  0 & 0 & 0
\end{array}
\right]
\]
\end{resolution}

%------------------------------------------------------------------------------%
\begin{resolution}{Solução}
A segunda linha representa
\[
0 = 0
\]

\pause
\bigskip
Portanto, o sistema possui \emph{infinitas} soluções
\end{resolution}

%------------------------------------------------------------------------------%
\endinput

%------------------------------------------------------------------------------%
\begin{question}
Use escalonamento para analisar o sistema
\[
  \systeme{
     x + y = 4,
    2x - y = 1
  }
\]
\end{question}

%------------------------------------------------------------------------------%
\begin{resolution}{Escalonamento}
\begin{columns}
\column{0.4\linewidth}
  \[\;\;\,
    \left[
    \begin{array}{rr|r}
      1 &  1 & 4 \\
      2 & -1 & 1
    \end{array}
    \right]
  \]
  \pause
\column{0.6\linewidth}
  Eliminar os elementos abaixo do primeiro pivô
  \\ \pause
  \(
    L_2 \leftarrow L_2 - 2 L_1
  \)
\end{columns}
\vspace{-\baselineskip}
\[
  \pause
  L'_2
  \;=\;
  L_2 - 2 L_1
  \pause
  \;=\;
  \left[
  \begin{array}{rr|r}
    2 & -1 & 1
  \end{array}
  \right]
  - 2
  \left[
  \begin{array}{rr|r}
    1 & 1 & 4
  \end{array}
  \right]
  \pause
  \;=\;
  \left[
  \begin{array}{rr|r}
    0 & -3 & -7
  \end{array}
  \right]
\]

\pause
\[
\left[
\begin{array}{rr|r}
  1 &  1 & 4 \\
  0 & -3 & -7
\end{array}
\right]
\]

\pause
O escalonamento produz dois pivôs e, portanto, uma única solução
\end{resolution}

%------------------------------------------------------------------------------%
\endinput


%------------------------------------------------------------------------------%
%\section{Lista Mínima}
%
%%------------------------------------------------------------------------------%
%\begin{frame}{Lista Mínima}
%  \begin{enumerate}
%    \item
%  \end{enumerate}
%
%  \vfill
%  Atenção: A prova é baseada no livro, não nas apresentações
%\end{frame}

%------------------------------------------------------------------------------%
\end{document}
%------------------------------------------------------------------------------%


